\nivelbasico


\begin{ejercicio}
Una partícula se mueve sobre el eje \(x\). En el instante

\[ t_i=\SI{2.0}{\second} \]

se encuentra en:

\[ x_i=\SI{3.0}{\meter}. \]

En el instante:

\[ t_f=\SI{8.0}{\second} \]

se encuentra en:

\[ x_f=\SI{27.0}{\meter}. \]

Determine:

\begin{enumerate}[label=\alph*)]
    \item el desplazamiento;
    \item el intervalo de tiempo;
    \item la velocidad media.
\end{enumerate}

\textit{Pista: la velocidad media relaciona el desplazamiento con el intervalo de tiempo.}
\end{ejercicio}

\begin{respuesta}
\[
\Delta x=\SI{24.0}{\meter},
\qquad
\Delta t=\SI{6.0}{\second},
\qquad
\bar{v}_x=\SI{4.0}{\meter\per\second}.
\]
\end{respuesta}

\begin{explicacion}
El desplazamiento es:

\[ \Delta x=x_f-x_i. \]

Sustituyendo:

\[
\Delta x
=
\SI{27.0}{\meter}
-
\SI{3.0}{\meter}
=
\SI{24.0}{\meter}.
\]

Por tanto:

\[ \boxed{\Delta x=\SI{24.0}{\meter}}. \]

El intervalo de tiempo es:

\[ \Delta t=t_f-t_i. \]

Entonces:

\[
\Delta t
=
\SI{8.0}{\second}
-
\SI{2.0}{\second}
=
\SI{6.0}{\second}.
\]

La velocidad media se define mediante:

\[ \bar{v}_x=\frac{\Delta x}{\Delta t}. \]

Sustituyendo:

\[
\bar{v}_x
=
\frac{\SI{24.0}{\meter}}
{\SI{6.0}{\second}}
=
\SI{4.0}{\meter\per\second}.
\]

Así:

\[ \boxed{\bar{v}_x=\SI{4.0}{\meter\per\second}}. \]

El signo positivo indica que el desplazamiento neto ocurrió en el sentido positivo del eje \(x\).
\end{explicacion}


\begin{ejercicio}
Una persona camina \(\SI{60}{\meter}\) hacia el este y después regresa \(\SI{20}{\meter}\) hacia el oeste. Todo el recorrido dura \(\SI{20}{\second}\).

Tome el este como sentido positivo y determine:

\begin{enumerate}[label=\alph*)]
    \item la distancia total recorrida;
    \item el desplazamiento;
    \item la rapidez media;
    \item la velocidad media.
\end{enumerate}

\textit{Pista: la rapidez media utiliza distancia; la velocidad media utiliza desplazamiento.}
\end{ejercicio}

\begin{respuesta}
\[
d=\SI{80}{\meter},
\qquad
\Delta x=+\SI{40}{\meter},
\]

\[
\bar{v}_{\mathrm{rap}}=\SI{4.0}{\meter\per\second},
\qquad
\bar{v}_x=+\SI{2.0}{\meter\per\second}.
\]
\end{respuesta}

\begin{explicacion}
La distancia recorrida es la suma de las longitudes de ambos tramos:

\[
d
=
\SI{60}{\meter}
+
\SI{20}{\meter}.
\]

Por tanto:

\[ \boxed{d=\SI{80}{\meter}}. \]

Para determinar el desplazamiento podemos tomar como origen el punto de partida.

Después del primer tramo:

\[ x=+\SI{60}{\meter}. \]

Al regresar \(\SI{20}{\meter}\):

\[
x_f
=
\SI{60}{\meter}
-
\SI{20}{\meter}
=
\SI{40}{\meter}.
\]

Entonces:

\[ \boxed{\Delta x=+\SI{40}{\meter}}. \]

La rapidez media es:

\[ \bar{v}_{\mathrm{rap}}=\frac{d}{\Delta t}. \]

Por tanto:

\[
\bar{v}_{\mathrm{rap}}
=
\frac{\SI{80}{\meter}}
{\SI{20}{\second}}
=
\boxed{\SI{4.0}{\meter\per\second}}.
\]

La velocidad media es:

\[ \bar{v}_x=\frac{\Delta x}{\Delta t}. \]

Así:

\[
\bar{v}_x
=
\frac{\SI{40}{\meter}}
{\SI{20}{\second}}
=
\boxed{+\SI{2.0}{\meter\per\second}}.
\]

La rapidez media es mayor porque tiene en cuenta toda la longitud recorrida, mientras que la velocidad media depende del cambio neto de posición.
\end{explicacion}


\begin{ejercicio}
Un móvil realiza movimiento rectilíneo uniforme. En \(t=0\) se encuentra en:

\[ x_0=-\SI{5.0}{\meter} \]

y se mueve con velocidad constante:

\[ v_x=+\SI{3.0}{\meter\per\second}. \]

Determine su posición en:

\[ t=\SI{4.0}{\second}. \]

\textit{Pista: para MRU, \(x=x_0+v_xt\).}
\end{ejercicio}

\begin{respuesta}
\[ \boxed{x=\SI{7.0}{\meter}}. \]
\end{respuesta}

\begin{explicacion}
Como la velocidad es constante, el movimiento es rectilíneo uniforme.

La ecuación de posición es:

\[ x=x_0+v_xt. \]

Sustituyendo:

\[
x
=
-\SI{5.0}{\meter}
+
(\SI{3.0}{\meter\per\second})
(\SI{4.0}{\second}).
\]

Entonces:

\[
x
=
-\SI{5.0}{\meter}
+
\SI{12.0}{\meter}.
\]

Por tanto:

\[ \boxed{x=\SI{7.0}{\meter}}. \]

El móvil comenzó a la izquierda del origen y, al desplazarse hacia el sentido positivo, terminó a la derecha de éste.
\end{explicacion}


\begin{ejercicio}
La posición de una partícula está descrita mediante:

\[ x(t)=12.0-2.50t, \]

donde \(x\) se expresa en metros y \(t\) en segundos.

Determine:

\begin{enumerate}[label=\alph*)]
    \item la posición inicial;
    \item la velocidad;
    \item la posición en \(t=\SI{4.0}{\second}\);
    \item el instante en el que atraviesa el origen.
\end{enumerate}

\textit{Pista: compare la ecuación con \(x=x_0+v_xt\).}
\end{ejercicio}

\begin{respuesta}
\[
x_0=\SI{12.0}{\meter},
\qquad
v_x=-\SI{2.50}{\meter\per\second},
\]

\[
x(\SI{4.0}{\second})=\SI{2.0}{\meter},
\qquad
t=\SI{4.80}{\second}.
\]
\end{respuesta}

\begin{explicacion}
La ecuación general del MRU es:

\[ x=x_0+v_xt. \]

Comparando con:

\[ x(t)=12.0-2.50t, \]

identificamos:

\[ \boxed{x_0=\SI{12.0}{\meter}} \]

y:

\[ \boxed{v_x=-\SI{2.50}{\meter\per\second}}. \]

El signo negativo indica que la partícula se mueve hacia el sentido negativo del eje.

Para \(t=\SI{4.0}{\second}\):

\[
x
=
12.0
-
2.50(4.0).
\]

Entonces:

\[ \boxed{x=\SI{2.0}{\meter}}. \]

Para encontrar cuándo atraviesa el origen hacemos:

\[ x=0. \]

Así:

\[
0=12.0-2.50t.
\]

Despejando:

\[
t=\frac{12.0}{2.50}.
\]

Por tanto:

\[ \boxed{t=\SI{4.80}{\second}}. \]
\end{explicacion}


\begin{ejercicio}
La siguiente gráfica representa la posición de un móvil en función del tiempo:

\begin{center}
\begin{tikzpicture}
\begin{axis}[
    width=0.78\textwidth,
    height=6cm,
    axis lines=left,
    xlabel={\(t\) (\si{\second})},
    ylabel={\(x\) (\si{\meter})},
    xmin=0,
    xmax=6,
    ymin=0,
    ymax=12,
    xtick={0,1,2,3,4,5,6},
    ytick={0,2,4,6,8,10,12},
    grid=both
]
\addplot[
    very thick,
    AzulPrincipal,
    domain=0:6,
    samples=2
]
{2+1.5*x};
\end{axis}
\end{tikzpicture}
\end{center}

Determine:

\begin{enumerate}[label=\alph*)]
    \item la posición inicial;
    \item la posición en \(t=\SI{4}{\second}\);
    \item la velocidad del móvil;
    \item si el movimiento es uniforme.
\end{enumerate}

\textit{Pista: en una gráfica \(x-t\), la pendiente representa la velocidad.}
\end{ejercicio}

\begin{respuesta}
\[
x_0=\SI{2.0}{\meter},
\qquad
x(\SI{4}{\second})=\SI{8.0}{\meter},
\]

\[
v_x=\SI{1.5}{\meter\per\second}.
\]

El movimiento es uniforme.
\end{respuesta}

\begin{explicacion}
La intersección de la recta con el eje vertical corresponde a la posición cuando:

\[ t=0. \]

De la gráfica:

\[ \boxed{x_0=\SI{2.0}{\meter}}. \]

En:

\[ t=\SI{4}{\second}, \]

la gráfica indica:

\[ \boxed{x=\SI{8.0}{\meter}}. \]

La velocidad corresponde a la pendiente:

\[ v_x=\frac{\Delta x}{\Delta t}. \]

Tomando dos puntos de la recta:

\[
(t_1,x_1)
=
(\SI{0}{\second},\SI{2}{\meter}),
\]

\[
(t_2,x_2)
=
(\SI{4}{\second},\SI{8}{\meter}),
\]

obtenemos:

\[
v_x
=
\frac{
\SI{8}{\meter}-\SI{2}{\meter}
}{
\SI{4}{\second}-0
}.
\]

Por tanto:

\[ \boxed{v_x=\SI{1.5}{\meter\per\second}}. \]

Como la gráfica es una recta de pendiente constante, la velocidad es constante y el movimiento es rectilíneo uniforme.
\end{explicacion}


\begin{ejercicio}
La posición de una partícula está dada por:

\[ x(t)=2.0+3.0t^2, \]

donde \(x\) está en metros y \(t\) en segundos.

Determine la velocidad instantánea en:

\begin{enumerate}[label=\alph*)]
    \item \(t=\SI{1.0}{\second}\);
    \item \(t=\SI{3.0}{\second}\).
\end{enumerate}

\textit{Pista: la velocidad instantánea se obtiene mediante \(v_x=\deriv{x}{t}\).}
\end{ejercicio}

\begin{respuesta}
\[
v_x(\SI{1.0}{\second})=\SI{6.0}{\meter\per\second},
\]

\[
v_x(\SI{3.0}{\second})=\SI{18.0}{\meter\per\second}.
\]
\end{respuesta}

\begin{explicacion}
La velocidad instantánea es la derivada de la posición respecto del tiempo:

\[ v_x=\deriv{x}{t}. \]

La posición es:

\[ x(t)=2.0+3.0t^2. \]

Derivando:

\[ v_x(t)=6.0t. \]

En:

\[ t=\SI{1.0}{\second}, \]

tenemos:

\[
v_x=6.0(1.0).
\]

Por tanto:

\[ \boxed{v_x=\SI{6.0}{\meter\per\second}}. \]

En:

\[ t=\SI{3.0}{\second}, \]

obtenemos:

\[
v_x=6.0(3.0).
\]

Así:

\[ \boxed{v_x=\SI{18.0}{\meter\per\second}}. \]

La velocidad no es constante. Por ello, este movimiento no es MRU.

La derivada indica la pendiente instantánea de la gráfica posición-tiempo.
\end{explicacion}


\nivelintermedio


\begin{ejercicio}
Un móvil se desplaza sobre una carretera recta de acuerdo con el siguiente recorrido:

\begin{itemize}
    \item entre \(t=0\) y \(t=\SI{10}{\second}\), pasa de \(x=0\) a \(x=\SI{30}{\meter}\);
    \item permanece en \(x=\SI{30}{\meter}\) durante \(\SI{5}{\second}\);
    \item después se mueve hasta \(x=\SI{10}{\meter}\), llegando allí en \(t=\SI{25}{\second}\).
\end{itemize}

Determine para todo el movimiento:

\begin{enumerate}[label=\alph*)]
    \item el desplazamiento;
    \item la distancia recorrida;
    \item la velocidad media;
    \item la rapidez media.
\end{enumerate}

\textit{Pista: el tiempo total incluye también el intervalo durante el cual el móvil permanece en reposo.}
\end{ejercicio}

\begin{respuesta}
\[
\Delta x=\SI{10}{\meter},
\qquad
d=\SI{50}{\meter},
\]

\[
\bar{v}_x=\SI{0.40}{\meter\per\second},
\qquad
\bar{v}_{\mathrm{rap}}=\SI{2.0}{\meter\per\second}.
\]
\end{respuesta}

\begin{explicacion}
El móvil comienza en:

\[ x_i=0 \]

y termina en:

\[ x_f=\SI{10}{\meter}. \]

Por tanto:

\[
\Delta x
=
x_f-x_i
=
\SI{10}{\meter}.
\]

Así:

\[ \boxed{\Delta x=\SI{10}{\meter}}. \]

La distancia del primer tramo es:

\[ d_1=\SI{30}{\meter}. \]

Durante el reposo:

\[ d_2=0. \]

En el último tramo se mueve desde:

\[ x=\SI{30}{\meter} \]

hasta:

\[ x=\SI{10}{\meter}, \]

por lo que:

\[ d_3=\SI{20}{\meter}. \]

Entonces:

\[
d
=
30+0+20
=
\boxed{\SI{50}{\meter}}.
\]

El intervalo total es:

\[ \Delta t=\SI{25}{\second}. \]

La velocidad media es:

\[
\bar{v}_x
=
\frac{\Delta x}{\Delta t}
=
\frac{\SI{10}{\meter}}
{\SI{25}{\second}}.
\]

Por tanto:

\[ \boxed{\bar{v}_x=\SI{0.40}{\meter\per\second}}. \]

La rapidez media es:

\[
\bar{v}_{\mathrm{rap}}
=
\frac{d}{\Delta t}
=
\frac{\SI{50}{\meter}}
{\SI{25}{\second}}.
\]

Así:

\[ \boxed{\bar{v}_{\mathrm{rap}}=\SI{2.0}{\meter\per\second}}. \]

La gran diferencia entre ambas magnitudes se debe al cambio de sentido del movimiento.
\end{explicacion}


\begin{ejercicio}
Dos automóviles \(A\) y \(B\) se mueven sobre la misma carretera recta.

En \(t=0\):

\[
x_{A0}=0,
\qquad
x_{B0}=\SI{100}{\meter}.
\]

Sus velocidades constantes son:

\[
v_A=+\SI{4.0}{\meter\per\second},
\qquad
v_B=-\SI{6.0}{\meter\per\second}.
\]

Determine:

\begin{enumerate}[label=\alph*)]
    \item el instante en que se encuentran;
    \item la posición del encuentro.
\end{enumerate}

\textit{Pista: en el encuentro ambos móviles tienen la misma posición en el mismo instante.}
\end{ejercicio}

\begin{respuesta}
\[
\boxed{t=\SI{10.0}{\second}},
\qquad
\boxed{x=\SI{40.0}{\meter}}.
\]
\end{respuesta}

\begin{explicacion}
Como ambos realizan MRU:

\[
x_A=x_{A0}+v_At,
\]

\[
x_B=x_{B0}+v_Bt.
\]

Para el móvil \(A\):

\[ x_A=4.0t. \]

Para \(B\):

\[ x_B=100-6.0t. \]

En el encuentro:

\[ x_A=x_B. \]

Por tanto:

\[
4.0t=100-6.0t.
\]

Agrupando:

\[ 10.0t=100. \]

Entonces:

\[ \boxed{t=\SI{10.0}{\second}}. \]

Para determinar la posición sustituimos en cualquiera de las ecuaciones:

\[
x_A
=
(\SI{4.0}{\meter\per\second})
(\SI{10.0}{\second}).
\]

Por tanto:

\[ \boxed{x=\SI{40.0}{\meter}}. \]

Comprobación con \(B\):

\[
x_B
=
100-6(10)
=
\SI{40}{\meter}.
\]

Ambos resultados coinciden, como debe ocurrir en el punto de encuentro.
\end{explicacion}


\begin{ejercicio}
Dos ciclistas se desplazan en la misma dirección sobre una carretera recta.

En \(t=0\):

\[
x_{A0}=0,
\qquad
x_{B0}=\SI{30}{\meter}.
\]

Sus velocidades son constantes:

\[
v_A=\SI{5.0}{\meter\per\second},
\qquad
v_B=\SI{2.0}{\meter\per\second}.
\]

Determine:

\begin{enumerate}[label=\alph*)]
    \item cuánto tiempo tarda \(A\) en alcanzar a \(B\);
    \item en qué posición ocurre el alcance.
\end{enumerate}

\textit{Pista: aunque se muevan en el mismo sentido, en el instante del alcance sus posiciones son iguales.}
\end{ejercicio}

\begin{respuesta}
\[
\boxed{t=\SI{10.0}{\second}},
\qquad
\boxed{x=\SI{50.0}{\meter}}.
\]
\end{respuesta}

\begin{explicacion}
Las ecuaciones de posición son:

\[
x_A=5.0t,
\]

\[
x_B=30+2.0t.
\]

Cuando \(A\) alcanza a \(B\):

\[ x_A=x_B. \]

Entonces:

\[
5.0t=30+2.0t.
\]

Por tanto:

\[
3.0t=30.
\]

Así:

\[ \boxed{t=\SI{10.0}{\second}}. \]

La posición del alcance es:

\[
x_A
=
(\SI{5.0}{\meter\per\second})
(\SI{10.0}{\second}).
\]

Entonces:

\[ \boxed{x=\SI{50.0}{\meter}}. \]

El ciclista \(A\) puede alcanzar a \(B\) porque su velocidad es mayor.
\end{explicacion}


\begin{ejercicio}
La gráfica siguiente representa la posición de un móvil:

\begin{center}
\begin{tikzpicture}
\begin{axis}[
    width=0.82\textwidth,
    height=6.5cm,
    axis lines=left,
    xlabel={\(t\) (\si{\second})},
    ylabel={\(x\) (\si{\meter})},
    xmin=0,
    xmax=10,
    ymin=-5,
    ymax=25,
    xtick={0,2,4,6,8,10},
    ytick={-5,0,5,10,15,20,25},
    grid=both
]

\addplot[
    very thick,
    AzulPrincipal
]
coordinates {
    (0,0)
    (4,20)
    (6,20)
    (10,0)
};

\end{axis}
\end{tikzpicture}
\end{center}

Determine:

\begin{enumerate}[label=\alph*)]
    \item la velocidad entre \(0\) y \(\SI{4}{\second}\);
    \item el estado de movimiento entre \(\SI{4}{\second}\) y \(\SI{6}{\second}\);
    \item la velocidad entre \(\SI{6}{\second}\) y \(\SI{10}{\second}\);
    \item el desplazamiento total;
    \item la distancia total recorrida.
\end{enumerate}

\textit{Pista: analice por separado la pendiente de cada tramo de la gráfica.}
\end{ejercicio}

\begin{respuesta}
\[
v_1=\SI{5.0}{\meter\per\second},
\qquad
v_2=0,
\qquad
v_3=-\SI{5.0}{\meter\per\second},
\]

\[
\Delta x=0,
\qquad
d=\SI{40}{\meter}.
\]
\end{respuesta}

\begin{explicacion}
En una gráfica posición-tiempo, la pendiente representa la velocidad.

Para el primer tramo:

\[
v_1
=
\frac{20-0}{4-0}
=
\boxed{\SI{5.0}{\meter\per\second}}.
\]

Entre:

\[
t=\SI{4}{\second}
\]

y:

\[
t=\SI{6}{\second},
\]

la posición permanece constante:

\[ x=\SI{20}{\meter}. \]

Por tanto:

\[ \boxed{v_2=0}. \]

El móvil está en reposo durante ese intervalo.

En el último tramo:

\[
v_3
=
\frac{0-20}{10-6}
=
-\SI{5.0}{\meter\per\second}.
\]

Así:

\[ \boxed{v_3=-\SI{5.0}{\meter\per\second}}. \]

El móvil comienza en:

\[ x_i=0 \]

y termina en:

\[ x_f=0. \]

Por tanto:

\[ \boxed{\Delta x=0}. \]

La distancia recorrida es:

\[
d
=
\SI{20}{\meter}
+
0
+
\SI{20}{\meter}.
\]

Entonces:

\[ \boxed{d=\SI{40}{\meter}}. \]

El móvil regresa al punto inicial, por lo que su desplazamiento es cero aunque haya recorrido una distancia no nula.
\end{explicacion}


\begin{ejercicio}
La velocidad de un móvil está representada en la siguiente gráfica:

\begin{center}
\begin{tikzpicture}
\begin{axis}[
    width=0.82\textwidth,
    height=6cm,
    axis lines=middle,
    xlabel={\(t\) (\si{\second})},
    ylabel={\(v_x\) (\si{\meter\per\second})},
    xmin=0,
    xmax=10.5,
    ymin=-3,
    ymax=4,
    xtick={0,2,4,6,8,10},
    ytick={-2,0,3},
    grid=both
]

\addplot[
    very thick,
    VerdePrincipal
]
coordinates {
    (0,3)
    (4,3)
};

\addplot[
    very thick,
    VerdePrincipal
]
coordinates {
    (4,-2)
    (10,-2)
};

\end{axis}
\end{tikzpicture}
\end{center}

El cambio instantáneo de velocidad en \(t=\SI{4}{\second}\) se idealiza en la gráfica.

Determine:

\begin{enumerate}[label=\alph*)]
    \item el desplazamiento entre \(0\) y \(\SI{4}{\second}\);
    \item el desplazamiento entre \(\SI{4}{\second}\) y \(\SI{10}{\second}\);
    \item el desplazamiento total;
    \item la distancia total;
    \item la velocidad media;
    \item la rapidez media.
\end{enumerate}

\textit{Pista: el área algebraica bajo una gráfica \(v_x-t\) representa desplazamiento.}
\end{ejercicio}

\begin{respuesta}
\[
\Delta x_1=\SI{12}{\meter},
\qquad
\Delta x_2=-\SI{12}{\meter},
\]

\[
\Delta x=0,
\qquad
d=\SI{24}{\meter},
\]

\[
\bar{v}_x=0,
\qquad
\bar{v}_{\mathrm{rap}}=\SI{2.4}{\meter\per\second}.
\]
\end{respuesta}

\begin{explicacion}
Entre \(0\) y \(\SI{4}{\second}\):

\[
v_x=\SI{3}{\meter\per\second}.
\]

El desplazamiento corresponde al área rectangular:

\[
\Delta x_1
=
v_x\Delta t
=
(\SI{3}{\meter\per\second})
(\SI{4}{\second}).
\]

Entonces:

\[ \boxed{\Delta x_1=\SI{12}{\meter}}. \]

Entre \(\SI{4}{\second}\) y \(\SI{10}{\second}\):

\[
v_x=-\SI{2}{\meter\per\second}.
\]

El intervalo es:

\[ \Delta t=\SI{6}{\second}. \]

Por tanto:

\[
\Delta x_2
=
(-\SI{2}{\meter\per\second})
(\SI{6}{\second})
=
-\SI{12}{\meter}.
\]

Así:

\[ \boxed{\Delta x_2=-\SI{12}{\meter}}. \]

El desplazamiento total es:

\[
\Delta x
=
12-12
=
\boxed{0}.
\]

La distancia total suma las magnitudes de ambos recorridos:

\[
d
=
12+12
=
\boxed{\SI{24}{\meter}}.
\]

La velocidad media es:

\[
\bar{v}_x
=
\frac{\Delta x}{\Delta t}
=
0.
\]

La rapidez media es:

\[
\bar{v}_{\mathrm{rap}}
=
\frac{\SI{24}{\meter}}
{\SI{10}{\second}}
=
\boxed{\SI{2.4}{\meter\per\second}}.
\]
\end{explicacion}


\begin{ejercicio}
La posición de una partícula está dada por:

\[ x(t)=5.0+4.0t-t^2, \]

donde \(x\) se mide en metros y \(t\) en segundos.

Determine:

\begin{enumerate}[label=\alph*)]
    \item la función velocidad \(v_x(t)\);
    \item la velocidad en \(t=\SI{1.0}{\second}\);
    \item la velocidad en \(t=\SI{3.0}{\second}\);
    \item el instante en que la velocidad es cero;
    \item la posición en ese instante.
\end{enumerate}

\textit{Pista: derive la función de posición y después analice el signo de la velocidad.}
\end{ejercicio}

\begin{respuesta}
\[
v_x(t)=4.0-2t,
\]

\[
v_x(\SI{1.0}{\second})=\SI{2.0}{\meter\per\second},
\]

\[
v_x(\SI{3.0}{\second})=-\SI{2.0}{\meter\per\second},
\]

\[
t=\SI{2.0}{\second},
\qquad
x=\SI{9.0}{\meter}.
\]
\end{respuesta}

\begin{explicacion}
La velocidad instantánea se obtiene mediante:

\[ v_x=\deriv{x}{t}. \]

Para:

\[ x(t)=5.0+4.0t-t^2, \]

derivamos término a término:

\[ \boxed{v_x(t)=4.0-2t}. \]

En:

\[ t=\SI{1.0}{\second}, \]

\[
v_x
=
4.0-2(1.0)
=
\boxed{\SI{2.0}{\meter\per\second}}.
\]

En:

\[ t=\SI{3.0}{\second}, \]

\[
v_x
=
4.0-2(3.0)
=
\boxed{-\SI{2.0}{\meter\per\second}}.
\]

Para encontrar cuándo la velocidad es cero:

\[
0=4.0-2t.
\]

Entonces:

\[ \boxed{t=\SI{2.0}{\second}}. \]

La posición en ese instante es:

\[
x(2)
=
5.0+4.0(2)-(2)^2.
\]

Así:

\[
x
=
5+8-4
=
\boxed{\SI{9.0}{\meter}}.
\]

Antes de \(t=\SI{2}{\second}\), la velocidad es positiva. Después de ese instante, es negativa.

Por tanto, en \(x=\SI{9}{\meter}\) la partícula cambia el sentido de su movimiento.
\end{explicacion}


\nivelavanzado


\begin{ejercicio}
Dos móviles \(A\) y \(B\) realizan movimientos rectilíneos uniformes descritos mediante:

\[
x_A=x_{A0}+v_At,
\]

\[
x_B=x_{B0}+v_Bt.
\]

Deduzca una expresión general para el instante de encuentro y analice:

\begin{enumerate}[label=\alph*)]
    \item qué condición debe cumplirse para que exista un instante de encuentro único;
    \item qué ocurre si \(v_A=v_B\) y \(x_{A0}\neq x_{B0}\);
    \item qué ocurre si \(v_A=v_B\) y \(x_{A0}=x_{B0}\).
\end{enumerate}
\end{ejercicio}

\begin{respuesta}
\[ \boxed{t_{\mathrm{enc}}=\frac{x_{B0}-x_{A0}}{v_A-v_B}}. \]

Para un encuentro único debe cumplirse:

\[ v_A\neq v_B \]

y el tiempo obtenido debe pertenecer al intervalo físico considerado.

Si las velocidades son iguales y las posiciones iniciales diferentes, nunca se encuentran.

Si las velocidades y posiciones iniciales son iguales, permanecen juntos en todo instante.
\end{respuesta}

\begin{explicacion}
En un encuentro los dos móviles ocupan la misma posición en el mismo instante:

\[ x_A=x_B. \]

Por tanto:

\[
x_{A0}+v_At
=
x_{B0}+v_Bt.
\]

Agrupamos los términos que contienen \(t\):

\[
(v_A-v_B)t
=
x_{B0}-x_{A0}.
\]

Si:

\[ v_A-v_B\neq0, \]

podemos dividir y obtener:

\[ \boxed{t_{\mathrm{enc}}=\frac{x_{B0}-x_{A0}}{v_A-v_B}}. \]

Sin embargo, obtener matemáticamente un valor de \(t\) no basta. Debe comprobarse que pertenezca al intervalo físico estudiado.

Si:

\[ v_A=v_B \]

pero:

\[ x_{A0}\neq x_{B0}, \]

la separación entre ellos permanece constante. No existe encuentro.

Si:

\[
v_A=v_B
\]

y además:

\[
x_{A0}=x_{B0},
\]

las dos ecuaciones de posición son idénticas. Los móviles coinciden en todo instante.

Este análisis muestra que una ecuación general también debe interpretarse físicamente y examinarse en sus casos especiales.
\end{explicacion}


\begin{ejercicio}
Dos móviles siguen las ecuaciones:

\[
x_A(t)=-10+4t,
\]

\[
x_B(t)=20-2t,
\]

con posiciones en metros y tiempo en segundos.

Sin comenzar directamente con una sustitución numérica:

\begin{enumerate}[label=\alph*)]
    \item interprete la posición inicial y el sentido de movimiento de cada móvil;
    \item determine si pueden encontrarse para \(t\geq0\);
    \item calcule el instante y la posición de encuentro;
    \item describa cómo se representaría el encuentro en una gráfica \(x-t\).
\end{enumerate}
\end{ejercicio}

\begin{respuesta}
\[
x_{A0}=-\SI{10}{\meter},
\qquad
v_A=+\SI{4}{\meter\per\second},
\]

\[
x_{B0}=+\SI{20}{\meter},
\qquad
v_B=-\SI{2}{\meter\per\second}.
\]

Se encuentran en:

\[
\boxed{t=\SI{5.0}{\second}},
\qquad
\boxed{x=\SI{10.0}{\meter}}.
\]

En una gráfica \(x-t\), las dos rectas se intersectan en ese punto.
\end{respuesta}

\begin{explicacion}
Para \(A\):

\[ x_A=-10+4t. \]

Su posición inicial es:

\[ x_{A0}=-\SI{10}{\meter}, \]

y su velocidad:

\[ v_A=+\SI{4}{\meter\per\second}. \]

Por tanto, se mueve hacia posiciones crecientes.

Para \(B\):

\[ x_B=20-2t. \]

Su posición inicial es:

\[ x_{B0}=+\SI{20}{\meter}, \]

y:

\[ v_B=-\SI{2}{\meter\per\second}. \]

Por tanto, se mueve hacia posiciones decrecientes.

Los dos móviles se aproximan, por lo que físicamente esperamos un encuentro.

En ese instante:

\[ x_A=x_B. \]

Entonces:

\[
-10+4t=20-2t.
\]

De donde:

\[
6t=30.
\]

Así:

\[ \boxed{t=\SI{5.0}{\second}}. \]

Sustituyendo en la ecuación de \(A\):

\[
x_A
=
-10+4(5)
=
\SI{10}{\meter}.
\]

Por tanto:

\[ \boxed{x=\SI{10}{\meter}}. \]

En una gráfica posición-tiempo, cada MRU aparece como una recta.

La recta de \(A\) tiene pendiente positiva y la de \(B\) pendiente negativa. El punto donde ambas rectas se cruzan representa el mismo tiempo y la misma posición para ambos cuerpos:

\[
(t,x)
=
(\SI{5.0}{\second},\SI{10}{\meter}).
\]
\end{explicacion}


\begin{ejercicio}
Un móvil realiza un recorrido rectilíneo durante \(\SI{30}{\second}\).

Durante los primeros \(\SI{10}{\second}\) se mueve con velocidad constante:

\[ v_1=+\SI{4.0}{\meter\per\second}. \]

Durante los siguientes \(\SI{8}{\second}\) permanece en reposo.

En el tiempo restante se mueve con velocidad constante:

\[ v_3=-\SI{5.0}{\meter\per\second}. \]

Si comienza en:

\[ x_0=-\SI{10}{\meter}, \]

determine:

\begin{enumerate}[label=\alph*)]
    \item la posición al finalizar cada etapa;
    \item la posición final;
    \item el desplazamiento total;
    \item la distancia total;
    \item la velocidad media;
    \item la rapidez media;
    \item una descripción cualitativa de la gráfica \(x-t\).
\end{enumerate}
\end{ejercicio}

\begin{respuesta}
\[
x_1=\SI{30}{\meter},
\qquad
x_2=\SI{30}{\meter},
\qquad
x_f=-\SI{30}{\meter},
\]

\[
\Delta x=-\SI{20}{\meter},
\qquad
d=\SI{100}{\meter},
\]

\[
\bar{v}_x=-\SI{0.667}{\meter\per\second},
\qquad
\bar{v}_{\mathrm{rap}}=\SI{3.33}{\meter\per\second}.
\]
\end{respuesta}

\begin{explicacion}
La posición inicial es:

\[ x_0=-\SI{10}{\meter}. \]

Durante los primeros \(\SI{10}{\second}\):

\[
\Delta x_1
=
v_1\Delta t_1
=
(\SI{4.0}{\meter\per\second})
(\SI{10}{\second})
=
\SI{40}{\meter}.
\]

Por tanto:

\[
x_1
=
-\SI{10}{\meter}
+
\SI{40}{\meter}
=
\boxed{\SI{30}{\meter}}.
\]

Durante los siguientes \(\SI{8}{\second}\):

\[ v_2=0. \]

Por tanto:

\[
\Delta x_2=0
\]

y:

\[ \boxed{x_2=\SI{30}{\meter}}. \]

El tiempo restante es:

\[
\Delta t_3
=
30-10-8
=
\SI{12}{\second}.
\]

En esta etapa:

\[
\Delta x_3
=
(-\SI{5.0}{\meter\per\second})
(\SI{12}{\second})
=
-\SI{60}{\meter}.
\]

La posición final es:

\[
x_f
=
\SI{30}{\meter}
-
\SI{60}{\meter}.
\]

Así:

\[ \boxed{x_f=-\SI{30}{\meter}}. \]

El desplazamiento total es:

\[
\Delta x
=
x_f-x_0
=
-30-(-10).
\]

Por tanto:

\[ \boxed{\Delta x=-\SI{20}{\meter}}. \]

La distancia recorrida es:

\[
d
=
|\Delta x_1|
+
|\Delta x_2|
+
|\Delta x_3|.
\]

Entonces:

\[
d
=
40+0+60
=
\boxed{\SI{100}{\meter}}.
\]

La velocidad media es:

\[
\bar{v}_x
=
\frac{-\SI{20}{\meter}}
{\SI{30}{\second}}
\approx
\boxed{-\SI{0.667}{\meter\per\second}}.
\]

La rapidez media es:

\[
\bar{v}_{\mathrm{rap}}
=
\frac{\SI{100}{\meter}}
{\SI{30}{\second}}
\approx
\boxed{\SI{3.33}{\meter\per\second}}.
\]

La gráfica \(x-t\) tendría:

\begin{itemize}
    \item un primer tramo recto con pendiente positiva;
    \item un tramo horizontal durante el reposo;
    \item un último tramo recto con pendiente negativa y de mayor magnitud que la primera.
\end{itemize}
\end{explicacion}


\begin{ejercicio}
Un automóvil realiza MRU. Se sabe que:

\[
x(\SI{2.0}{\second})=\SI{12}{\meter},
\]

y:

\[
x(\SI{8.0}{\second})=-\SI{6}{\meter}.
\]

Determine:

\begin{enumerate}[label=\alph*)]
    \item la velocidad del automóvil;
    \item su posición en \(t=0\);
    \item la ecuación \(x(t)\);
    \item el instante en que atraviesa el origen;
    \item el sentido de movimiento.
\end{enumerate}
\end{ejercicio}

\begin{respuesta}
\[
v_x=-\SI{3.0}{\meter\per\second},
\]

\[
x_0=\SI{18}{\meter},
\]

\[
\boxed{x(t)=18-3t},
\]

\[
\boxed{t=\SI{6.0}{\second}}.
\]

El móvil se desplaza en el sentido negativo del eje.
\end{respuesta}

\begin{explicacion}
Como el movimiento es uniforme, la velocidad puede obtenerse a partir de dos posiciones:

\[
v_x
=
\frac{x_2-x_1}{t_2-t_1}.
\]

Sustituyendo:

\[
v_x
=
\frac{
-\SI{6}{\meter}
-
\SI{12}{\meter}
}{
\SI{8}{\second}
-
\SI{2}{\second}
}.
\]

Entonces:

\[
v_x
=
\frac{-\SI{18}{\meter}}
{\SI{6}{\second}}
=
\boxed{-\SI{3.0}{\meter\per\second}}.
\]

Ahora utilizamos:

\[ x=x_0+v_xt. \]

Con el dato:

\[
x(\SI{2}{\second})=\SI{12}{\meter},
\]

tenemos:

\[
12=x_0-3(2).
\]

Por tanto:

\[
x_0=18.
\]

Así:

\[ \boxed{x_0=\SI{18}{\meter}}. \]

La ecuación de posición es:

\[ \boxed{x(t)=18-3t}. \]

Para encontrar cuándo atraviesa el origen:

\[ x=0. \]

Entonces:

\[
0=18-3t.
\]

De donde:

\[ \boxed{t=\SI{6.0}{\second}}. \]

La velocidad es negativa durante todo el movimiento, por lo que el automóvil se desplaza en el sentido negativo del eje \(x\).
\end{explicacion}


\begin{ejercicio}
Un estudiante afirma:

\begin{quote}
``Si la velocidad media de un móvil durante cierto intervalo es cero, entonces necesariamente permaneció en reposo durante todo ese intervalo''.
\end{quote}

Determine si la afirmación es correcta. Construya un ejemplo numérico que justifique su respuesta y calcule tanto la velocidad media como la rapidez media.
\end{ejercicio}

\begin{respuesta}
La afirmación es falsa.

Por ejemplo, un móvil puede recorrer \(\SI{20}{\meter}\) hacia la derecha y regresar \(\SI{20}{\meter}\) al punto inicial en \(\SI{10}{\second}\).

En ese caso:

\[
\bar{v}_x=0,
\]

pero:

\[
\bar{v}_{\mathrm{rap}}=\SI{4.0}{\meter\per\second}.
\]
\end{respuesta}

\begin{explicacion}
Tomemos como ejemplo un móvil que comienza en:

\[ x_i=0. \]

Se mueve hasta:

\[ x=+\SI{20}{\meter}, \]

y posteriormente vuelve al origen:

\[ x_f=0. \]

Supongamos que todo el movimiento dura:

\[ \Delta t=\SI{10}{\second}. \]

El desplazamiento total es:

\[
\Delta x
=
x_f-x_i
=
0.
\]

Por tanto:

\[
\bar{v}_x
=
\frac{\Delta x}{\Delta t}
=
0.
\]

Así:

\[ \boxed{\bar{v}_x=0}. \]

Sin embargo, la distancia recorrida es:

\[
d
=
20+20
=
\SI{40}{\meter}.
\]

La rapidez media es:

\[
\bar{v}_{\mathrm{rap}}
=
\frac{\SI{40}{\meter}}
{\SI{10}{\second}}
=
\boxed{\SI{4.0}{\meter\per\second}}.
\]

El móvil se desplazó durante el intervalo, pero terminó en su posición inicial.

Por ello, una velocidad media igual a cero no implica reposo.
\end{explicacion}


\begin{ejercicio}
Demuestre que, para cualquier movimiento rectilíneo durante un intervalo de tiempo \(\Delta t>0\), la rapidez media nunca puede ser menor que la magnitud de la velocidad media:

\[ \bar{v}_{\mathrm{rap}}\geq|\bar{v}_x|. \]

Indique cuándo se cumple la igualdad.
\end{ejercicio}

\begin{respuesta}
\[ \boxed{\bar{v}_{\mathrm{rap}}\geq|\bar{v}_x|}. \]

La igualdad se cumple cuando:

\[ d=|\Delta x|, \]

es decir, cuando el móvil no cambia de sentido durante el intervalo.
\end{respuesta}

\begin{explicacion}
Por definición:

\[ \bar{v}_{\mathrm{rap}}=\frac{d}{\Delta t}. \]

La magnitud de la velocidad media es:

\[
|\bar{v}_x|
=
\left|
\frac{\Delta x}{\Delta t}
\right|.
\]

Como:

\[ \Delta t>0, \]

podemos escribir:

\[
|\bar{v}_x|
=
\frac{|\Delta x|}{\Delta t}.
\]

Para cualquier movimiento rectilíneo sabemos que:

\[ d\geq|\Delta x|. \]

Dividiendo ambos lados entre el intervalo positivo \(\Delta t\):

\[
\frac{d}{\Delta t}
\geq
\frac{|\Delta x|}{\Delta t}.
\]

Por tanto:

\[ \boxed{\bar{v}_{\mathrm{rap}}\geq|\bar{v}_x|}. \]

La igualdad se cumple cuando:

\[ d=|\Delta x|. \]

Esto sucede si el móvil se desplaza sin cambiar de sentido durante el intervalo considerado.

Si existe al menos un cambio de sentido que produzca recorrido adicional:

\[ d>|\Delta x|, \]

y por ello:

\[ \bar{v}_{\mathrm{rap}}>|\bar{v}_x|. \]

La desigualdad no es simplemente una propiedad algebraica: expresa la diferencia física entre considerar todo el camino recorrido y considerar solamente el cambio neto de posición.
\end{explicacion}